\chapter{Le fonds Grothendieck de Montpellier}\label{ch:montpellier}

\section{La source}

Le fonds Alexandre Grothendieck a été donné à l'université de Montpellier par Jean
Malgoire et catalogué en 2015--2016 par Hélène Rodriguez et Frédéric Troilo, sous la
direction de Sophie Dikoff. Il compte environ 28\,000 pages ; environ 18\,000 peuvent être
communiquées, la correspondance des tiers ne pouvant l'être sans autorisation. La partie
ouverte en ligne (\url{https://grothendieck.umontpellier.fr/archives-grothendieck/}) compte
178 dossiers, environ 16\,000 pages, de 1949 à 1991. Chaque dossier est un seul PDF, de 4 à
270~Mo ; l'ensemble pèse environ 28~Go. Tout est en français.

C'est la source sur laquelle la méthode a été inventée, et la plus difficile du livre,
techniquement et juridiquement.

\section{Ce qui a coincé}

\paragraph{Le serveur refuse d'être affiché ailleurs.} Trois constats, faits à
l'automne 2026 :
\begin{itemize}
  \item il envoie \texttt{X-Frame-Options: SAMEORIGIN} : aucun autre site ne peut afficher
        ses PDF dans une page ;
  \item son certificat TLS a expiré le 10 décembre 2025, et un cadre intégré reste blanc,
        sans message ;
  \item il n'envoie aucun en-tête CORS : une page d'un autre site ne peut pas lire les
        fichiers.
\end{itemize}

\paragraph{La réponse : un relais.} Le projet a écrit un petit serveur Node
(\fichier{relay/server.mjs}), hébergé gratuitement, qui va chercher le PDF chez Montpellier
et le renvoie au lecteur. Il faut Node parce qu'il faut accepter le certificat expiré, ce
qu'un service sans serveur ne permet pas. Le relais transmet les requêtes partielles
(\emph{range requests}) : ouvrir la page 221 d'un volume de 204~Mo ne déplace que 256~Ko.
Il n'accepte que \texttt{GET} et \texttt{HEAD} sur une adresse par cote, ne transmet aucun
cookie, ne garde rien. L'hébergement gratuit s'endort au bout de quinze minutes : la
première page met une minute à venir. Rien de cela n'est satisfaisant ; c'est ce qu'il
fallait pour qu'un lecteur voie le fac-similé à côté de la transcription sans que le
projet stocke une seule image.

\paragraph{La copie de travail.} Pour transcrire, on extrait un lot de vingt pages dans un
dossier ignoré par \texttt{git}. Le premier outil essayé (poppler) faisait d'un lot de
vingt pages un fichier de 42~Mo, plus lourd que le volume source ; pdf-lib en fait 7~Mo.
Ensuite, chaque page est découpée en six tuiles qui se chevauchent, à 700~dpi. Cet outil a
été écrit après un lot de seize pages qui avait demandé quatre-vingt-dix recadrages à la
main.

\paragraph{Les pages.} Grothendieck écrivait souvent au dos de documents sans rapport : ces
pages-là sont sautées, sans être décrites. Certaines ont été numérisées à l'envers. Trois
numérotations se superposent (celle du PDF, décalée d'une unité par une page de garde,
celle de l'archive, celle de l'auteur). Presque rien n'est daté ; les datations entre
crochets de l'inventaire sont recopiées telles quelles.

\section{Les droits, et ce qu'on a décidé}

Grothendieck est mort le 13 novembre 2014. Ses droits patrimoniaux courent jusqu'à la fin
de 2084 et appartiennent à ses héritiers ; son droit moral, dont le droit de divulgation,
est perpétuel. Une transcription reproduit le texte : elle est couverte par ces droits,
qu'elle soit tapée par une personne ou par une machine.

Le projet a d'abord mis tout le dépôt sous CC0, en écrivant que le fonds « n'appartenait à
personne ». C'était faux ; c'est retiré, et le fichier \fichier{RIGHTS.md} dit maintenant :
\begin{itemize}
  \item le code, les scripts, le site et les preuves Lean sont sous CC0 ;
  \item les transcriptions ne sont sous aucune licence du projet, qui n'en revendique aucun
        droit et n'autorise personne à les reproduire ;
  \item chaque page porte le filigrane « Édition de démonstration » ;
  \item aucun fac-similé n'est stocké ni redistribué ;
  \item si l'université ou les ayants droit demandent le retrait, il est fait, sans
        discussion ni délai ; un interrupteur de déploiement retire du site tout le texte
        des transcriptions sans rien effacer ;
  \item aucun dépôt sur Zenodo ni identifiant pérenne (DOI), parce qu'un identifiant
        pérenne rendrait le retrait impossible ; il en irait de même pour un modèle
        entraîné sur ces transcriptions.
\end{itemize}
La question de ce que permet la mention « libre accès » du fonds a été posée par écrit au
service des archives et n'a pas reçu de réponse. En attendant, le projet lit cette mention
comme couvrant une édition de démonstration, filigranée et retirable.

Les lettres ont leur propre règle : celles des tiers ne sont pas diffusées, les lettres
publiées avec Serre restent résumées, et les personnes privées ne sont pas nommées. Le
registre des lettres compte 93 entrées.

\section{Ce qui a été produit}

À la fin d'octobre 2026 : 638 lots transcrits, couvrant 164 dossiers et environ 9\,400
pages, chacun avec sa lecture modernisée ; une seconde passe sur tuiles pour 61 lots ; des
énoncés démontrés en Lean pour 73 dossiers ; des vérifications par une personne, ponctuelles
(125 lectures dans 54 dossiers, 31 lettres, 57 dates) ; cinq articles. Aucun lot n'est
vérifié en entier par une personne.

\begin{exemple}
La leçon de Montpellier est que l'obstacle technique se contourne en une semaine, et que
l'obstacle juridique ne se contourne pas. Le relais est un bricolage ; le filigrane, la
promesse de retrait et l'absence de DOI sont des décisions, prises tôt et écrites, qui
permettent au projet d'exister tant que personne ne s'y oppose.
\end{exemple}
